\Page{9}{\textenglish{02 / FORECAST}}{מה התמונה באמת מספרת}{45}{52}

\Lead{בתמונת רחוב אחת יש מעיל, תיק, צבעים ושאלה מי נכנסה בכלל לפריים. בשביל לבדוק אם מאפיין חוזר מספיק פעמים, מישהו צריך קודם להגדיר אותו. האם המערכת סופרת מעילים ארוכים או את כל הבגדים שמגיעים עד אמצע השוק? ראייה ממוחשבת יכולה לסווג מאפיינים בתמונות רבות. ההגדרה וההרכב של המדגם קובעים איך לקרוא את המספרים. \SourceRef{5}\SourceRef{6}}

\vspace{5mm}

{\fontsize{14.8}{23.1}\selectfont

\StoryColumns{\textbf{\textenglish{WGSN}} מציגה שיטה שמחברת תצוגות, רשתות, חיפוש, נתוני היצע בחנויות ועמדות של צרכניות. החברה מתארת זיהוי תמונות ומודלים של למידת מכונה לצד מומחי תחום. עבור חזאית, התועלת היא האפשרות לבדוק אם כמה סוגי אותות מתיישבים זה עם זה. צבע שעלה על המסלול יכול להישאר בגדר הצעה אמנותית. אם במקביל הוא מופיע במבחר של חנויות ומעורר חיפושים, נפתחת שאלה אחרת על שלב האימוץ שלו. \SourceRef{4}}{\textbf{\textenglish{Heuritech}} מדגישה זיהוי לבוש במדגמי תמונות מרשתות חברתיות וחיזוי השכיחות של מאפיינים לאורך זמן. לפי החברה, אפשר להבחין בין הופעת פריט בקרב מאמצות מוקדמות לבין הופעתו בקבוצות רחבות יותר. אותו גוון יכול לעלות בשוק אחד בלי לקבל אותה תשומת לב בשוק אחר. הפלטפורמה נותנת למותג דרך לראות את הפער, והצוות צריך לבדוק מה קורה אצלו במסחר. \SourceRef{5}}

\par}

\vspace{6mm}

\Panel{גם כשאותות עולים יחד, הם עשויים להתייחס לדברים שונים. תמונה מלמדת מה הופיע במדגם. חיפוש מעיד על עניין או סקרנות. רשומת מכירה מספרת מה נרכש במקום ובמחיר מסוימים, כשהמלאי היה זמין. לכן כדאי לשאול איזה נתון עומד מאחורי התחזית לפני שמשתמשים בה להגדלת הזמנה. פרסומי שתי החברות מתארים שיטות וכלים. הם אינם מספקים מבחן השוואתי על אותם נתונים. \SourceRef{4}\SourceRef{5}}
